\subsection{受控复现实验}
\label{sec:results-round2}

既有文献留下三个决定上述结论可外推程度的问题：有监督模型是否也会停在随机水平；schema 结果是否只是某一个检查点的习惯；第二个推理模型是否显示同样的轴。本节用预先部署的对照回答这三个问题，并针对 schema 结果检验第四个问题——跨家族的一般性；其余跨家族主张仍未解决，并如实标注。

\subsubsection{监督基线并不处于随机水平}
\label{sec:results-baseline}

我们用 ImageNet 初始化训练 DenseNet-121~\citep{huang2017densely}：$224$\,px、类别加权 BCE、AdamW、固定八个 epoch，不做早停、不做基于测试集的选择。它只见过 3,871 张 \texttt{internal\_train} 图像，与 1,506 例评测研究没有任何患者重叠。在同样的七个病理标签上，它取得\textbf{宏 AUROC 0.733}（95\% 患者聚类自助法区间 0.718--0.749，725 名患者，1,000 次重采样；逐标签 0.620--0.805；表~\ref{tab:baseline}），在其自身 1,352 例患者不重叠的 \texttt{internal\_test} 划分上为 0.724（VLM 评测使用该划分中的 1,054 例子集；第~\ref{sec:protocol} 节）。

\begin{table}[t]
\centering
\small
\caption{干净 NIH 七个病理标签的宏 AUROC（$n=1{,}506$）。基线只用 3,871 张无患者重叠的图像训练；六个 VLM 即表~\ref{tab:clean} 的矩阵。}
\label{tab:baseline}
\begin{tabular}{lrr}
\toprule
模型 & 宏 AUROC & 家族 \\
\midrule
\textbf{DenseNet-121（监督）} & \textbf{0.7329} & CNN，本文训练 \\
MedGemma-4B & 0.6270 & 医学 VLM \\
InternVL3.5-8B & 0.6066 & 通用 VLM \\
MedGemma-1.5 & 0.5852 & 医学 VLM \\
Qwen3-VL-8B & 0.5455 & 通用 VLM \\
Qwen2.5-VL-7B & 0.5013 & 通用 VLM \\
CheXagent-8B & 0.5001 & 胸片专用 \\
\bottomrule
\end{tabular}
\end{table}

因此，这些标签从这些图像上是可学的，水平 0.73，而训练集只有 3,871 张——不到公开训练划分 86,524 张的 5\%。``没有可用判别力''说的是模型，不是任务或单纯的标签噪声。基线也给出反方向的界：0.73 并非临床级，故本文的负结果针对的是\emph{测量可识别性}，而不是胸片分类本身不可能。

\subsubsection{schema 依赖性在第二个医学模型上复现，但映射不同}
\label{sec:results-schema2}

我们在 MedGemma-1.5 上重做五套 schema 扫描（15,060 次调用）。核心结果复现：在空白图像上，\textbf{五套 schema 的病理召回全部为 0.000}，而真实图像上为 0.493--0.614（表~\ref{tab:r2schema}）。没有任何措辞能制造图像依赖。

\emph{未}复现的是每套 schema 落入哪种失效形态。MedGemma-4B 有两种（自信判正常，或被迫弃权）；MedGemma-1.5 出现三种：自信判正常（v0、v1、v3 的空白回答中 \texttt{no\_finding} 阳性率 100\%）、\emph{契约坍缩}（v2 仍给标签但在 99.9\% 的空白回答中丢掉概率对象）、以及拒绝断言（v4 在 100\% 的空白回答中判 \texttt{no\_finding} 阴性，并同样契约坍缩）。真实图像召回最高的是 v4（0.614，比 v0 高 6.3\,\pp{}），而它的空白召回仍为 0.000：提示词移动的是工作点，不是模型获取证据的能力。

正确的表述因此比第~\ref{sec:results-schema} 节的``契约决定捷径的形式''更窄：契约与模型\emph{共同}决定形式，而且这一联合映射不可跨检查点迁移。可迁移的是对基准真正重要的不变性：真实图像召回与空白图像行为由不同因素决定，因此任何 schema 都不能只靠其中之一来验证。

\subsubsection{空白图像结果跨模型家族成立；契约签名不成立}
\label{sec:results-xfamily}

为区分 MedGemma 习惯与契约效应，我们在 Qwen3-VL-8B（不同家族的通用模型）上重做扫描（表~\ref{tab:r2xfamily}）。空白图像结果精确复现：三套 schema 下空白图像的病理召回均为 \textbf{0.000}，真实图像为 0.258--0.306。可观察签名没有复现：Qwen3 的 v0 是自信正常（空白回答中 \texttt{no\_finding} 阳性 100\%），但八个请求概率只给出一个，概率完整度空白 12.5\%、真实 47.1\%；v1 保持同样的自信正常答案，同时把完整度恢复到 100\% 与 95.2\%；v4 拒绝正常断言（空白中 \texttt{no\_finding} 阳性 0\%），并在 98\% 的真实回答中契约坍缩。两个模型出现相同的三种失效形态，但 variant 到形态的映射不同。因此契约覆盖是 schema 的属性，而不是模型的属性；空白图像依赖则同时是二者的属性。

\begin{table}[t]
\centering
\footnotesize
\setlength{\tabcolsep}{3.5pt}
\caption{Qwen3-VL-8B 上的跨家族 schema 扫描（每格 $n=1{,}506$）。三套 schema 下空白图像的病理召回均为 0.000，因此空白依赖不是 MedGemma 的习惯。可观察签名不同：v0 自信正常但只给一个概率，v1 自信正常且契约完整，v4 拒绝正常断言并契约坍缩。}
\label{tab:r2xfamily}
\begin{tabular}{lrrrrrr}
\toprule
方案 & 真实召回 & 空白召回 & 空白 \texttt{no\_finding} 阳性 & 违约（真实/空白） & 完整度（真实/空白） \\
\midrule
v0 简单 & 0.306 & \textbf{0.000} & 1.000 & 0.611 / 1.000 & 0.471 / 0.125 \\
v1 逐标签二分类 & 0.258 & \textbf{0.000} & 1.000 & 0.054 / 0.000 & 0.952 / 1.000 \\
v4 正常性施压 & 0.263 & \textbf{0.000} & 0.000 & 0.983 / 1.000 & 0.849 / 0.875 \\
\bottomrule
\end{tabular}
\end{table}

\subsubsection{措辞不是机制，但它决定空白图像上的答案}
\label{sec:results-paraphrase}

我们在 MedGemma-4B 上对 v0 契约做了两个改写版本（标签集合、状态规则、JSON 结构完全一致），对真实与空白图像各跑一遍（6,024 次调用）。真实图像上，发现级答案可复现：与 v0 基线的一致率为 90.9\% 与 91.4\%，Cohen's $\kappa$ 为 0.624 与 0.642。schema 结论不是单一措辞的产物。

空白图像上，两个改写版本与基线只在一个标签上分歧。v0 在 100\% 的空白回答中判 \texttt{no\_finding} 阳性；p2 相同（一致率 1.000）；p1 在 100\% 中判阴性（一致率 0.000）。病理标签在三者之间完全一致（全阴性，一致率 1.000）。也就是说，空白图像上的答案由措辞决定、不含任何图像信息——这正是为什么零模型前提需要一个空白对照\emph{并且}一个措辞对照，而不是二者之一。

该对照也在 Qwen3-VL 上运行（表~\ref{tab:r2para}）。两个改写版本都复现真实图像答案（p1 一致率 93.1\%，$\kappa$ 0.723；p2 90.3\%，$\kappa$ 0.701），但都不翻转空白图像答案：与 v0 的空白一致率均为 1.000，而 p1 会把 MedGemma-4B 的空白答案完全翻转。因此措辞设定的是模型特定的先验：同一改写把 MedGemma 的空白答案变负，却让 Qwen3 的自信正常答案保持不变。

\subsubsection{第二个推理模型：同一工作点，新的失效模式}
\label{sec:results-second}

Qwen3-VL-8B-Thinking 具有原生推理通道（chat 模板预填 \texttt{<think>}，模型在回答前闭合它），因此可察觉性轴第一次能在第二个模型上读出。我们在全部 1,506 例研究上跑了七个格子（clean 加六个 severe）。只有闭合推理块\emph{并且}其后能解析出 JSON 的回复才算可用答案，其余作为实测结果保留（表~\ref{tab:r2recon}）。

\textbf{覆盖率本身是可测量，且随退化下降。} 在 5,120 token 预算下，10,542 个格子中有 9,381 个给出可用答案（89.0\%）；对 40 个格子的 8,192 token 抽样再补 20 个（最终 89.2\%）。覆盖率为 clean 92.6\%、分辨率 severe 80.9\%、亮度 severe 84.3\%（表~\ref{tab:r2think}）。失败并非``加大预算即可消除的截断''：仍然失败的回复中位长度在 2,560 上限时为 9.8k 字符、3,584 时为 13.5k、5,120 时为 19.3k，在 8,192 的 40 格随机抽样中为 28.1k——它们会撑满给它的任何预算；8,192 抽样只能救回 50\%（表~\ref{tab:r2budget}）。不终止是模型行为，我们按行为报告。这是 overthinking 文献在视觉语言模型上的对应：超过某个点，追加推理 token 的边际价值为负，并可能放弃模型本来已经得到的答案~\citep{wu2026whenmore,zhou2026overthinking}；被截断的推理链可能比不推理更糟~\citep{su2026broken}。这里的失效更严重——回复根本无法解析——因此预算曲线是该模型的测量属性，而不只是成本。

\begin{table}[t]
\centering
\footnotesize
\setlength{\tabcolsep}{3.5pt}
\caption{第二个推理模型，severe 格（每格 $n=1{,}506$）。覆盖率为 5,120 token 预算下给出可用答案的研究比例；$\Delta$ 召回为与 clean 在同一批共同作答研究上的配对差（noise $n=1{,}238$；resolution $n=1{,}132$）。AUROC 为七个病理标签的宏平均；阳性率为被判阳性的病理槽位比例。}
\label{tab:r2think}
\begin{tabular}{lrrrrr}
\toprule
格子 & 覆盖率 & 链长(k) & 阳性率 & 宏 AUROC & $\Delta$ 召回 \\
\midrule
clean & 0.926 & 8.73 & 0.250 & 0.547 & --- \\
blur & 0.932 & 8.84 & 0.189 & 0.547 & $-5.5$\,\pp{} \\
brightness & 0.843 & 9.56 & 0.213 & 0.550 & $-5.1$\,\pp{} \\
contrast & 0.924 & 8.99 & 0.280 & 0.545 & $+2.2$\,\pp{} \\
JPEG & 0.923 & 9.37 & 0.240 & 0.537 & $-1.8$\,\pp{} \\
noise & 0.885 & 9.06 & 0.173 & 0.527 & $\mathbf{-10.3}$\,\pp{} \\
resolution & \textbf{0.809} & \textbf{10.22} & \textbf{0.125} & 0.549 & $\mathbf{-13.2}$\,\pp{} \\
\bottomrule
\end{tabular}
\end{table}

\textbf{工作点机制复现，方向相反。} 分辨率损失下阳性率由 0.250 降至 0.125，\texttt{no\_finding} 阳性率由 0.283 升至 0.450，而宏 AUROC 在各格中稳定在 0.527--0.550。由于覆盖率同时下降，我们在配对子集（两格都给出可用答案的研究）上重算该比率：0.250 $\rightarrow$ 0.128，即这一移动不是覆盖率伪影（附录~\ref{app:round2}，表~\ref{tab:r2posrate}）。模型并未丢失排序信息——它本来就没有——它把工作点移向``正常''。这与第~\ref{sec:results-op} 节互为镜像：那里干净阳性率高的模型在退化下上移、召回上升；这里干净阳性率仅 25\% 的模型下移、召回下降。二者都是工作点移动，都不是鲁棒性变化。

逐标签看，配对损失经聚类自助法与多重比较控制后依然显著：胸腔积液 $-8.8$ $[-11.9,-5.8]$\,\pp{}、肺炎 $-11.0$ $[-16.4,-5.5]$\,\pp{}（$q<0.001$）、心影增大 $-6.0$ $[-9.9,-1.9]$\,\pp{}（$q=0.014$）、肺不张 $-2.2$ 与肺水肿 $-4.2$\,\pp{}（$q=0.053$），实变与气胸不显著（表~\ref{tab:r2label}）。

\textbf{轴对通用，象限不通用。} 以覆盖率损失作为该模型的可察觉性读数、以配对召回变化作为损伤，两者再次可分离（与 $|\Delta|$ 的 Spearman 0.49，$p=0.33$；与带符号 $\Delta$ 的为 $-0.54$，$p=0.27$；$n=6$）。在同一 $|\Delta|$ 定义下，MedGemma-1.5 的标记--损伤相关也是 0.49，但六个点使两个估计都无法与零区分：我们主张同一\emph{方向}，而不是共同斜率。噪声换位了：对 MedGemma-1.5，噪声是静默退化（标记率 0.46\%，倒数第二低）而损伤与分辨率损失相当；对 Qwen3-VL-Thinking，噪声是响亮的（覆盖率损失 11.5\%，第三高），并造成 10.3\,\pp{} 的召回损失。\emph{静默失效是监测通道的属性，而不是退化的属性}：盯着推理是否终止的通道能看见噪声，盯着某个特定标记词的通道看不见。

这也解释了为什么第二个模型没有复现反转 3 的正号。该反转需要阳性率有上行空间；一个只把 25\% 槽位判阳的模型几乎没有空间，退化只能把它往下推。按第四前提报告工作点，正是使两个模型可比的原因。
